% Appendices E, F and G of the New In ML 2026 camera-ready paper: the composition proofs, the
% experiments in full, and reproducibility.

\section{Composition at a factor of two: proofs}\label{app:dichotomy}

This appendix proves the results of Section~\ref{sec:composition} and states those the body only
cites. Throughout, ``the class'' is the set of all $\ell_1$-non-expansive repairs defined there,
and ``the walk'' is our covering transplant of Algorithm~1 of \citet{dinitz2021}, defined in
Proposition~\ref{prop:walkbound}, with residuals fixed at their initial values.

\subsection{No additive constant for the class}

\begin{proof}[Proof of Theorem~\ref{thm:noadditiveclass}]
\emph{One block.} Every prediction in this proof is non-negative, so $z \in \dfeas$ is admissible exactly when $\|z - y\|_1 \le \|\hat y - y\|_1$ for all $y \in \dfeas$. Take $A = \{a\}$, $B = \{b_1, b_2, b_3\}$, $w = (1, 2, 2)$ and $\hat y = (0 \mid 0, 2, 2)$. With $y_a = t$ the cheapest completion has value
$t + (1-t)^+ + 2(2-t)^+$, which is $5 - 2t$ on $[0,1]$, $4 - t$ on $[1,2]$ and $t$ beyond. So
$\opt = 2$, attained only at $t = 2$ with the $B$-coordinates all $0$. The optimal face is
therefore the single point $y^* = (2 \mid 0, 0, 0)$, and $\hat y$ agrees with it at $b_1$ alone,
so $\|y^* - \hat y\|_0 = 3$. The point $y_1 = (0 \mid 1, 2, 2)$ is feasible, at $\ell_1$ distance
$1$ from $\hat y$, so any admissible $z$ has $\|z - y_1\|_1 \le 1$. Agreeing with $y^*$ at $a$,
$b_2$ or $b_3$ costs $2$. Agreeing at $b_1$ costs the whole budget and forces
$z = (0 \mid 0, 2, 2) = \hat y$, which is infeasible. So $z$ agrees with $y^*$ nowhere:
$\|y^* - z\|_0 = 4$.

\emph{Every block at once.} For a prediction $\hat y$ on a system with feasible set $\dfeas$ put
$e(z) = \sup_{y \in \dfeas} \bigl(\|z - y\|_1 - \|\hat y - y\|_1\bigr)$, so that $z$ is admissible
exactly when $e(z) \le 0$. On the block above, $e \ge 0$ everywhere: with $y_2 = (1 \mid 0, 2, 2)$,
also feasible at distance $1$ from $\hat y$, the triangle inequality gives
$\|z - y_1\|_1 + \|z - y_2\|_1 \ge \|y_1 - y_2\|_1 = 2$, so one of the two terms is at least
$1 = \|\hat y - y_i\|_1$. Now take $k$ disjoint copies of the block, with all cross weights $0$.
Cross constraints read $y_i + y_j \ge 0$ and are implied by $y \ge 0$, so $\dfeas$ is the product
of the blocks' feasible sets. The optimum $\opt$ is the sum of the blocks' optima, attained only
when every block is optimal, so the face is the single point $y^* = (y^*_1, \dots, y^*_k)$ and
$\|y^* - \hat y\|_0 = 3k$. Also $e(z) = \sum_g e_g(z_g)$, since the supremum of a sum of functions
of disjoint variables over a product set is the sum of the suprema. Since each $e_g \ge 0$,
$e(z) \le 0$ forces $e_g(z_g) \le 0$ in every block, in particular $\|z_g - y_1\|_1 \le 1$, and the
one-block argument gives $\|y^*_g - z_g\|_0 = 4$ in every block.
\end{proof}

In its smallest form the witness saturates the coordinate count of Section~\ref{sec:prelim},
$4k = |A|+|B|$. Corollary~\ref{cor:sparse} shows that the obstruction does not depend on
this.

\begin{proof}[Proof of Corollary~\ref{cor:sparse}]
Pad the witness. Adjoin $M$ further rows and $N$ further columns, all of whose weights are $0$,
and predict $0$ on them. All cross weights are $0$, so $\dfeas$ is still a product and $\opt$ is
still $2k$. The face is still a single point, now with a $0$ in each adjoined coordinate, which
$\hat y$ already matches, so $\|y^* - \hat y\|_0 = 3k$ as before. And $e$ still decomposes: the
summand of an adjoined coordinate is $e_i(z_i) = z_i$, so $e(z) \le 0$ forces $z_i = 0$ there.
Every admissible $z$ therefore still has $\|y^* - z\|_0 = 4k$, now with $|A| = k + M$ and
$|B| = 3k + N$. The increase does not depend on $M$ and $N$, which can therefore be arbitrarily large.
Taking $M = N = k$ leaves the increase $k$ at $4k$ against a coordinate count of $6k$ and against
$\|y^* - \hat y\|_0 + \min(|A|,|B|) = 5k$, while still attaining $\|y^* - \hat y\|_0 + \tau(G(p))$
exactly. Taking $M = N = k^2$ leaves the same increase $k$ from a prediction wrong on
$3k = \Theta(\sqrt{n})$ of $n = 2k^2 + 4k$ coordinates.
\end{proof}

\begin{remark}[Scope of the theorem]\label{rem:noadditivescope}
(i) The theorem excludes an additive constant but not a factor: the witness has ratio only $4/3$.
An exhaustive search over $1 \times 2$ and $1 \times 3$ systems with a unique optimal dual and
weights and predictions in $\{0, \dots, 4\}$ finds no larger ratio, and no witness on
$1 \times 2$. Theorem~\ref{thm:coverscan} gives an explicit factor~$2$, and
Theorem~\ref{thm:threehalves} fixes the class constant on bipartite covering instances at $3/2$.
(ii) The theorem assumes non-expansion toward \emph{every} feasible point, which
Theorem~\ref{thm:nonexp} gives and on which the obstruction rests. It does not speak to the weaker
property that the runtime bounds use, an $\ell_1$ bound toward optimal duals only, under which
Corollary~3.3 of \citet{chen2022} composes;
there additivity is possible (Proposition~\ref{prop:faceadditive}). Face-only non-expansion is
also unobstructed on this family, whose face is the single point $y^*$.
\end{remark}

\subsection{The constant of the class is \texorpdfstring{$3/2$}{3/2}}

Theorems~\ref{thm:noadditiveclass} and~\ref{thm:coverscan} put the class constant between $4/3$
and $2$. On bipartite covering instances it is exactly $3/2$:

\begin{theorem}[The constant of the class is $3/2$]\label{thm:threehalves}
Let $p = \max(\hat y, 0)$. As in Section~\ref{sec:composition}, $z \in \dfeas$ is
\emph{admissible} for $\hat y$ if $\|z - y\|_1 \le \|p - y\|_1$ for every $y \in \dfeas$, and a map
is an $\ell_1$-non-expansive repair exactly when every output is admissible. Any choice of an
admissible point per prediction is one (returning $\repair_\sigma$ elsewhere), so the admissible
points are exactly the outputs the class can return on $\hat y$. On bipartite covering instances:
\begin{enumerate}
\item[(a)] For every optimal dual $f$ and every maximum-weight matching $M$ there is an admissible
$z$ with
\begin{align*}
  \|f - z\|_0 \;&\le\; |\mathrm{supp}(p - f)| + \bigl|\{ v : p_v < f_v,\ p_{m(v)} > f_{m(v)} \}\bigr| \\
  \;&\le\; |\mathrm{supp}(p - f)| + \min\bigl(|C^-|, |C^+|\bigr),
\end{align*}
where $m(v)$ is the $M$-partner of $v$ and $C^-$, $C^+$ are the coordinates on which $p$ lies
below, respectively above, $f$. In particular, from a prediction at $\|y^* - \hat y\|_0 = r$ some
admissible $z$ has $\|y^* - z\|_0 \le r + \lfloor r/2 \rfloor \le \tfrac{3}{2}\,r$, and
$\|y^* - z\|_0 \le r$ whenever $p \le f$ or $p \ge f$ coordinatewise.
\item[(b)] For every $k \ge 1$ there is an instance with $|A| = k$, $|B| = 2k+1$, a unique optimal
dual, and a prediction with $\|y^* - \hat y\|_0 = 2k+1$ on which \emph{every} admissible $z$ has
$\|y^* - z\|_0 = 3k+1 = (2k+1) + \lfloor (2k+1)/2 \rfloor = |A| + |B|$: every admissible output
differs from the face on every coordinate.
\end{enumerate}
Consequently the best multiplicative $\ell_0$ constant over all $\ell_1$-non-expansive repairs, of
any running time, is exactly $3/2$: the map sending $\hat y$ to the point of (a) is a member of the
class with $\|y^* - \repair(\hat y)\|_0 \le \lfloor 3\|y^* - \hat y\|_0 / 2 \rfloor$ on every
bipartite covering instance, and by (b) no member has a smaller constant, the ratios $(3k+1)/(2k+1)$
increasing to $3/2$. The bound of (a) is attained for every odd $r$ (for $r \ge 3$ by (b), for
$r = 1$ by any feasible non-optimal prediction, whose only admissible output is $p$ itself).
The order in (a) is read off an
optimal dual and a maximum-weight matching; the theorem says nothing about repairs that do not
consult the optimum.
\end{theorem}

\begin{proof}
(a) Write $C = \mathrm{supp}(p - f) = C^- \sqcup C^+$. By complementary slackness ($M$ is an optimal solution of the matching LP, since the bipartite matching polytope is integral), every coordinate with $f_v > 0$, in particular every $v \in C^-$, is matched in $M$ by an edge of positive weight, and every edge of $M$ is tight at $f$. Starting from $p$, run two rounds of full-residual raises on currently violated constraints, each constraint fired at most once.
\emph{Round one:} for each $v \in C^-$ in turn, fire $(v, m(v))$ if it is violated. Since $M$ is a
matching, no two constraints of round one share an endpoint. So $v$ is still at $p_v$ when its
constraint is decided, and so is $m(v)$, unless $m(v)$ lies in $C^-$ and was processed earlier.
This is the only case in which a round-one constraint can meet an already-raised endpoint, and
it is handled below for either order. If $m(v)$ is correct ($p_{m(v)} = f_{m(v)}$), the residual
is exactly $f_v - p_v$, so $v$ reaches $f_v$ and $m(v)$ leaves $f$: one coordinate fixed, one
broken. If $m(v) \in C^-$, then whichever of the two is processed first, both endpoints end at or
above $f$ on those coordinates and stay wrong. Neither is counted as fixed. If $m(v) \in C^+$,
either nothing fires or $v$ ends strictly below $f_v$, while $m(v)$ stays above: no change. After
round one the coordinates still below $f$ form a subset $L$ of
$\{ v \in C^- : m(v) \in C^+ \}$, and since $m$ is injective $|L| \le \min(|C^-|, |C^+|)$.

\emph{Round two:} every constraint still violated has an endpoint in $L$, since
$z_i + z_j < w_{ij} \le f_i + f_j$ forces an endpoint below $f$, and coordinates only rise. For each $v \in L$ in turn, fire its violated constraint of largest residual, if any. This satisfies every constraint at $v$ from then on and changes at most one further coordinate.

After round two nothing is violated, so the result is feasible. As in part (a) of the proof of Theorem~\ref{thm:coverscan} below, it is $\repair_\sigma(\hat y)$ for $\sigma$ listing the fired constraints in firing order, hence admissible by Theorem~\ref{thm:nonexp}. It differs from $f$ on at most $|C| + |L|$ coordinates. After round one the wrong set is $C$ minus the fixed coordinates plus their broken partners, of size at most $|C|$. Each round-two firing adds at most one wrong coordinate, its partner, whether that partner was correct, already broken, or fixed in round one (a fixed coordinate lies in $C$).

Now take $f$ at Hamming distance
$r$ from $\hat y$. Then $|C| \le r$ (where $\hat y_i < 0$, either $f_i = 0 = p_i$ or $f_i > 0 \ne \hat y_i$)
and $\min(|C^-|,|C^+|) \le \lfloor |C|/2 \rfloor$. If $p \le f$ or $p \ge f$, then $L = \emptyset$.

(b) Take $A = \{a_1, \dots, a_k\}$, $B = \{b_1, \dots, b_k, c_1, \dots, c_{k+1}\}$, weights
$w_{a_i b_i} = 1$ and $w_{a_i c_i} = w_{a_i c_{i+1}} = 2$, all others $0$, and
$\hat y = (0, \dots, 0 \mid 0, \dots, 0, 2, \dots, 2)$: rows and $b$'s at $0$, $c$'s at $2$, so
$p = \hat y$. \emph{The face is the point $f$ with $f_{a_i} = 2$ and $f = 0$ elsewhere.} It is
feasible with value $2k$, and the matching $\{(a_i, c_i)\}$ has weight $2k$, so $\opt = 2k$. For any
optimal $y$, summing $y_{a_i} + y_{c_i} \ge 2$ over $i$ gives
$\sum_i y_{a_i} + \sum_{i \le k} y_{c_i} \ge 2k = \sum y$, so $y_{c_{k+1}} = 0$, $y_b = 0$, and each
$y_{a_i} + y_{c_i}$ equals $2$. Then $y_{c_{k+1}} = 0$ forces $y_{a_k} \ge 2$, hence $y_{c_k} = 0$,
hence $y_{a_{k-1}} \ge 2$, and so on down to $y_{c_1} = 0$ and $y_{a_i} = 2$ for all $i$. So
$\|y^* - \hat y\|_0 = 2k+1$: the $k$ rows and the $k+1$ columns $c_j$.
\emph{Every admissible $z$ differs from $f$ on all $3k+1$ coordinates.} Admissible points exist:
$\repair_\sigma(\hat y) = (1, \dots, 1 \mid 1, \dots, 1, 2, \dots, 2)$ for every order. For each
$x \in \prod_i \{a_i, b_i\}$ the point $y_x = \hat y + \sum_i e_{x_i}$, which raises $a_i$ or $b_i$
by one in every block, is feasible ($a_i + b_i \ge 1$, and every $c_j$ is $2$) and at $\ell_1$
distance $k$ from $\hat y$. Write $z_{a_i} = \alpha_i$, $z_{b_i} = \beta_i$, $z_{c_j} = 2 + \gamma_j$,
with $z \ge 0$. Admissibility against $y_x$ reads $\sum_i q_i(x_i) + \sum_j |\gamma_j| \le k$ with
$q_i(a_i) = |\alpha_i - 1| + \beta_i$ and $q_i(b_i) = \alpha_i + |\beta_i - 1|$; since the blocks are
chosen independently, the maximum over $x$ of these $2^k$ inequalities is
\[
  \sum_i \max\bigl(q_i(a_i),\ q_i(b_i)\bigr) + \sum_j |\gamma_j| \;\le\; k ,
\]
and each maximum is at least $1$: for $t \ge 0$, $|t - 1| + t \ge 1$ with equality exactly when
$t \le 1$, so $q_i(a_i) + q_i(b_i) \ge 2$. Hence $\gamma = 0$, so every $c_j$ stays at $2 \ne 0$. In
every block both quantities equal $1$. By the equality case this forces
$\alpha_i, \beta_i \in [0, 1]$ and then $(1 - \alpha_i) + \beta_i = 1$, i.e.\ $\alpha_i = \beta_i$,
while feasibility of $(a_i, b_i)$ forces $\alpha_i = \beta_i \ge 1/2$. So $z_{a_i} \le 1 \ne 2$ and
$z_{b_i} \ge 1/2 \ne 0$: all $3k+1$ coordinates differ from $f$. For $k = 1$ this is the gadget of
Theorem~\ref{thm:noadditiveclass}. For $k \ge 2$ the family chains $k$ copies of it through the
shared columns $c_2, \dots, c_k$. Each shared column constrains two rows at the cost of one wrong
coordinate while the face stays the single point $f$, and this sharing raises the ratio.
\end{proof}

The family of (b) saturates, $3k + 1 = |A| + |B|$, as the witness of
Theorem~\ref{thm:noadditiveclass} does, and it strengthens that theorem: the same increase
$k$ on $3k+1$ rather than $4k$ coordinates.

The upper bound as constructed consults an optimal dual and a maximum-weight matching, and
computing either is a full solve. The open question is whether any near-linear-time repair reaches
$3/2$, or even beats $2$. The cover-directed scan has factor~$2$, which
Proposition~\ref{prop:chainfam} shows is tight in the limit.

\subsection{The cover-directed scan}

\begin{remark}[The remaining parts of Theorem~\ref{thm:coverscan}]\label{rem:coverscanfull}
Let $S$ be any vertex cover of $G(p)$ and $\repair_S$ the cover-directed scan of
Theorem~\ref{thm:coverscan}. Part (a) is the claim of Theorem~\ref{thm:coverscan} that $\repair_S(\hat y)$ is feasible,
changes at most $2|S|$ coordinates of $p$, and equals $\repair_\sigma(\hat y)$ for an order
$\sigma$ that depends on $\hat y$. In addition:
\begin{enumerate}
\item[(b)] For every optimal $f$, $\mathrm{supp}(p - f)$ is a vertex cover of $G(p)$ contained in $\mathrm{supp}(\hat y - f)$; hence $\tau(G(p)) \le \|y^* - \hat y\|_0$.
\item[(c)] $\|y^* - \repair_S(\hat y)\|_0 \le \|y^* - \hat y\|_0 + |S|$.
\end{enumerate}
With $S$ the set of endpoints of any maximal matching of $G(p)$, found in $O(|\mathcal{C}|)$ time and needing no bipartiteness, $|S| \le 2\tau(G(p))$ and the bound is $\|y^* - \hat y\|_0 + 2\tau(G(p)) \le 3\|y^* - \hat y\|_0$. On a non-bipartite system the face distances are minima over the LP optimal face, which may be fractional (the integral-enumeration shortcut of Section~\ref{sec:prelim} needs bipartiteness).
\end{remark}

\begin{proof}[Proof of Theorem~\ref{thm:coverscan} and Remark~\ref{rem:coverscanfull}]
(a) Coordinates only rise, so a satisfied constraint stays satisfied. Firing the constraint of
largest residual $r$ at $i$ raises $y_i$ by $r$, which is at least the residual of every other
constraint at $i$, so afterwards every constraint at $i$ holds. Every edge of $G(p)$ has an
endpoint in $S$, and the constraints outside $G(p)$ held at $p$, so after the loop everything
holds. Each firing changes its two endpoints and there are at most $|S|$ firings. Let $\sigma$ list
the fired constraints in firing order, followed by all the others in any order: scanning $\sigma$
reproduces the same firings from the same states, and every remaining constraint is satisfied
when reached, so $\repair_\sigma(\hat y) = \repair_S(\hat y)$.

(b) If $(i,j)$ is violated at $p$ then $p_i + p_j < w_{ij} \le f_i + f_j$, so $p$ and $f$ differ
at $i$ or at $j$. Also, $\mathrm{supp}(p - f) \subseteq \mathrm{supp}(\hat y - f)$: where
$\hat y_i < 0$, either $f_i = 0 = p_i$ or $f_i > 0 \ne \hat y_i$. Taking $f$ at minimum Hamming
distance from $\hat y$ gives $\tau(G(p)) \le \|y^* - \hat y\|_0$.

(c) Let $C = \mathrm{supp}(p - f)$ with $f$ as in (b). A constraint fired by the scan is violated
when fired, hence violated at $p$, hence an edge of $G(p)$, so $C$ contains one of its endpoints:
of the two coordinates a firing changes, at most one lies outside $C$. The output differs from
$f$ at most on $C$ and on the changed coordinates outside $C$, so
$\|f - \repair_S(\hat y)\|_0 \le |C| + \#\mathrm{firings} \le \|y^* - \hat y\|_0 + |S|$, and the
minimum over the face is no larger. Finally, a minimum cover of a bipartite graph has size
$\nu(G(p))$ (K\H{o}nig). Either side of the bipartition is a cover, so
$\tau(G(p)) \le \min(|A|,|B|)$. The $2|M|$ endpoints of a maximal matching $M$ cover every edge,
and $|M| \le \tau(G(p))$.
\end{proof}

\begin{remark}[Reading the bound]\label{rem:coverscan}
First, the additive term is attained, and the factor~$2$ is tight in the limit for the scan as
stated. The chain family of Proposition~\ref{prop:chainfam} below ends at
$\|y^* - \hat y\|_0 + \tau(G(p))$ with ratio $(2k+1)/(k+1)$, under one of its two minimum covers,
in index order and with no ties. The worst case comes from the theorem's free choice of cover and
order. On $3{,}000$ random instances (shapes up to $4 \times 4$, weights in $\{0,\dots,5\}$,
predictions in $\{-2,\dots,6\}$), the worst ratio over every minimum cover, every processing order
of it and every tie among largest residuals is $4/3$ (a sample maximum, not a bound). The scan's
additive term is also at most Theorem~\ref{thm:sharpwarm}'s at every shape, since
$\min\{\|y^* - \hat y\|_0, |A|, |B|\} \le 2\min(|A|,|B|) - 1$, with equality only at
$\min(|A|,|B|) = 1$.

Second, the scan ends at face distance $1$ on both order-dependence families of Sections~\ref{sec:delicate}
and~\ref{sec:sorted}. On the family of Theorem~\ref{thm:order} the cover is $\{a\}$ and the largest
residual is at $b_{d-1}$, so it fires once and ends at face distance $1$ where the forward scan
ends at $d$. On the balanced family $\mathcal{B}_n$ of Definition~\ref{def:balancedfam} the cover
is $\{a_1\}$, the largest residual is at $b_n$, and it ends at $1$ where every weakly-descending
order ends at $n$ (Theorem~\ref{thm:balancedlb}). Both violated graphs are stars with a unique
largest residual. On them the walk also ends at face distance at most $1$, under every start and
tie-break, since it raises the centre by its largest remaining residual whether it starts there or
at a leaf. On a star the cover-directed scan fires
once, and there its selection rule is that of the classical maximal-residual relaxation method of
\citet{agmon1954} and \citet{motzkin1954}: a most violated constraint first. The full-residual
raise takes the place of the projection step.

Third, the cover determines where to raise; the raise itself is the usual
two-sided one. Raising only the cover vertices, each by its largest residual, is feasible, touches
at most $|S|$ coordinates (exactly $|S|$ for a minimum cover), and satisfies (c) of
Remark~\ref{rem:coverscanfull} with the same right-hand side. But it is one-sided and loses
non-expansion, like the walk. The two-sided full-residual raise keeps Theorem~\ref{thm:nonexp} at
no cost in the bound.
\end{remark}

\begin{proposition}[Theorem~\ref{thm:coverscan} is tight in the limit for the cover-directed scan]\label{prop:chainfam}
For $k \ge 2$ let $A_i = 2i + 1$ and take rows $a_1, \dots, a_k$, columns $b_1, \dots, b_k, c$ with
$w_{a_i b_i} = A_i$, $w_{a_{i+1} b_i} = A_{i+1} - 1$ (indices modulo $k$, so $w_{a_1 b_k} = A_1 - 1$),
$w_{a_i c} = A_i$, all other weights $0$, and $\hat y = 0$ except $\hat y_c = A_k$. The optimal face
is the single point with $y_{a_i} = A_i$ and $0$ elsewhere, so $\|y^* - \hat y\|_0 = k + 1$, and
$\tau(G(p)) = k$. The cover-directed scan run with the minimum cover $\{b_1, \dots, b_k\}$ in that
order meets no tie and ends at face distance $2k + 1 = \|y^* - \hat y\|_0 + \tau(G(p))$, a ratio
$(2k+1)/(k+1) \to 2$; on the same instance an admissible output at distance $k + 1$ exists.
\end{proposition}

\begin{proof}
The matching $\{(a_i, b_i)\}$ has weight $\sum_i A_i$ and the stated point is feasible with that
value, so it is optimal. For any optimal $y$, summing $y_{a_i} + y_{b_i} \ge A_i$ gives
$\sum_i (y_{a_i} + y_{b_i}) \ge \sum_i A_i = \sum y$, so $y_c = 0$ and every $y_{a_i} + y_{b_i} = A_i$,
and then $y_{a_i} + y_c \ge A_i$ forces $y_{a_i} = A_i$, $y_{b_i} = 0$. At $p$ the violated constraints
are the $2k$-cycle $a_1 b_1 a_2 b_2 \cdots a_k b_k a_1$, whose vertex cover number is $k$. At $b_i$
with $i < k$ the two violated constraints have residuals $A_i - z_{a_i} \le A_i$ and $A_{i+1} - 1$,
and the second is larger, so the scan fires $(a_{i+1}, b_i)$: $a_{i+1}$ ends at $A_{i+1} - 1$, one
short of its face value, and $b_i$ leaves $0$. At $b_k$ the residuals are $1$ and $A_1 - 1 = 2$, so
it fires $(a_1, b_k)$. Every row is one short, every $b_i$ and $c$ are off, and no residual was ever
tied. Firing instead the edge $(a_i, b_i)$ of each row (the row-side cover, or the tight edges
first) reaches the face point in every row, at distance $k + 1$. This output is
$\repair_\sigma(\hat y)$ for an order $\sigma$ listing those edges first, hence admissible by
Theorem~\ref{thm:nonexp}.
\end{proof}

The two constants therefore differ: the class can always reach $r + \lfloor r/2 \rfloor$, while an
explicit member can be forced to $r + \tau$. The difference comes from the optimal face, which the
construction of Theorem~\ref{thm:threehalves}(a) uses. An explicit repair loses in the choice
of which side of the violated graph to cover, and in what order: on the chain family the row side
adds no wrong coordinate and the column side ends at the largest distance the bound allows.

\subsection{The composition}

\begin{proof}[Proof of Theorem~\ref{thm:composed}]
\emph{Feasibility and cost.} Coordinates only decrease, and a packing constraint only gains slack
when its endpoints decrease, so no satisfied constraint is ever re-violated and every excess is
non-increasing throughout. When $i \in S$ is processed and its largest current excess is $e > 0$
at the constraint $(i,j)$, lowering $y_i$ and $y_j$ by $e$ lowers the excess of every other
constraint at $i$ by at least $e$, and each such excess was at most $e$. So after its turn every
constraint at $i$ holds for the rest of the scan. Every edge of $G(\hat y)$ has an endpoint in $S$, and
a constraint outside $G(\hat y)$ held at $\hat y$, so the output is feasible. It is integral
because $c$ and $\hat y$ are and each step subtracts an integral excess. There is at most one
firing per vertex of $S$, so at most $|S|$ firings and $O(m)$ work in the scan. A minimum vertex
cover of the bipartite graph $G(\hat y)$ is read off a maximum matching by K\H{o}nig's theorem in
$O(m\sqrt{n})$ time.

\emph{The count, against every optimal dual at once.} Fix any optimal dual $f$ and let
$C = \mathrm{supp}(\hat y - f)$. A constraint violated at $\hat y$ has
$\hat y_i + \hat y_j > c_{ij} \ge f_i + f_j$, so $\hat y$ and $f$ differ at $i$ or at $j$: $C$ is
a vertex cover of $G(\hat y)$, and hence $\tau \le |C| = \|f - \hat y\|_0$. Since excesses never
increase, every fired constraint was already violated at $\hat y$, hence is an edge of
$G(\hat y)$ and has an endpoint in $C$: of the two coordinates a firing changes, at most one lies
outside $C$. The output therefore differs from $f$ at most on $C$ and on the at most $|S|$ changed
coordinates outside $C$:
$\|f - \hat y'\|_0 \le |C| + |S| = \|f - \hat y\|_0 + \tau \le 2\,\|f - \hat y\|_0$.

\emph{The composition.} $\hat y'$ is a feasible integer dual, so \citet[Thm.~3.2]{chen2022}
applies to it and costs $O(m\sqrt{n} + (m + n\log n)\|y^* - \hat y'\|_0)$ against the optimal dual
$y^*$ their bound fixes. The count above holds for that $y^*$ in particular, and the $O(m\sqrt n)$
spent on the cover is absorbed by the first term.
\end{proof}

The count holds for every optimal dual, so the composition works against whichever one their
bound fixes. Reading $\|y^* - \hat y\|_0$ as the distance to the whole optimal face would be
strictly stronger, and it would need their bound to hold for every optimal dual at once, which
they do not state. The scan is the packing mirror of the cover-directed scan of
Theorem~\ref{thm:coverscan}, and the proof above mirrors that theorem's proof.
Proposition~\ref{prop:walkbound} proves the same count for the walk. For Algorithm~1 of
\citet{dinitz2021} itself, in the packing form in which it is stated (the repair
\citeauthor{chen2022}'s Corollary~3.3 names), we claim no $\ell_0$ bound.

\subsection{The walk}

\begin{proof}[Proof of Proposition~\ref{prop:walkbound}]
Feasibility is Lemma~5 of \citet{dinitz2021}: an edge $(i,j)$ is deleted when its first-visited endpoint $i$ is
raised, and then $\delta_i \ge r_{ij}$. For the count, fix $f$ at minimum Hamming distance from
$\hat y$ and let $C = \mathrm{supp}(p - f)$, by Remark~\ref{rem:coverscanfull}(b) a vertex cover of $G(p)$ with $|C| \le \|\hat y - f\|_0 = \|y^* - \hat y\|_0$. One pass of the inner loop visits vertices $i_1, \dots, i_L$, all
raised, each adjacent in $G(p)$ to the next at the time of the step, and stops at a vertex $j$
with no remaining edge. That vertex is not raised then, and never later, since edges are only ever
deleted. Distinct passes visit disjoint vertex sets. Consecutive raised vertices of a pass are
adjacent, so $C$ contains one of each pair, and at most $\lceil L/2 \rceil$ of the pass's raised
vertices lie outside $C$. The path $i_1, \dots, i_L, j$ has $L + 1$ vertices and so contains
$\lfloor (L+1)/2 \rfloor = \lceil L/2 \rceil$ pairwise disjoint edges. These paths are
vertex-disjoint across passes (a terminal has no edge left, so it is never visited or a terminal
again). So the union of those edges is a matching of $G(p)$, and
$\sum_{\mathrm{passes}} \lceil L/2 \rceil \le \nu(G(p)) = \tau(G(p))$. Hence the output differs
from $f$ on at most $|C| + \tau(G(p))$ coordinates. This is the path structure of their Lemma~6.
Their proof observes that the stepped edges form vertex-disjoint paths (they write ``paths and
cycles''; only paths arise, since visited vertices are deleted), and uses this, through a
half-integral matching, to bound the $\ell_1$ movement $\sum_i \delta_i$ within a factor~$2$ of
their correction LP. We use it to count coordinates; like theirs, the argument never consults the
tie-breaks.

For the family, $\opt = 5$, the face is the segment from $(2, 2 \mid 0, 0, 1, 0)$ to
$(3, 2 \mid 0, 0, 0, 0)$, and $\hat y$ agrees with it only at $b_1$. The two violated constraints
$(a_1, b_1)$ and $(a_2, b_1)$ both have residual $1$. Whichever vertex starts, the walk raises
$b_1$, alone or together with one row, by $1$, and breaks the one agreement: every output
differs from the face everywhere. Disjoint copies decouple as in the proof of
Theorem~\ref{thm:noadditiveclass}.
\end{proof}

The factor~$2$ is tight in the limit: an explicit \emph{ladder family} forces the walk arbitrarily
close to it.

\begin{proposition}[Proposition~\ref{prop:walkbound} is tight in the limit for the walk]\label{prop:ladderfam}
For $m \ge 2$ set $d_m = m + 1$, $d_i = d_{i+1} + i + 1$ for $i < m$, and
$s_i = d_i - i$ for $i < m$. Take rows $x_1, \dots, x_m$, columns $v_1, \dots, v_m$ with
\[
  w_{x_i v_i} = d_i, \qquad w_{x_i v_{i+1}} = d_{i+1}, \qquad w_{x_{i+1} v_i} = s_i
  \quad (i < m),
\]
all other weights $0$, and $\hat y = 0$, so $p = \hat y$. Then $\opt = \sum_i d_i$, the
optimal face is the polytope
\[
  F \;=\; \bigl\{ (\varepsilon \mid d - \varepsilon) \;:\;
  d_i \ge \varepsilon_i \ge \varepsilon_{i+1} \ge 0, \;\;
  \varepsilon_i - \varepsilon_{i+1} \le i \bigr\},
\]
$\|y^* - \hat y\|_0 = m$, and $G(p)$ consists of all $3m - 2$ positive-weight
constraints, with $\tau(G(p)) = m$. The walk started at $x_1$ meets no tie, traverses
$x_1 v_1 x_2 v_2 \cdots x_m v_m$ in a single pass, and ends at face distance exactly
$2m - 2 = \|y^* - \hat y\|_0 + \tau(G(p)) - 2$: a ratio
\[
  \frac{2m - 2}{m} \;\longrightarrow\; 2 .
\]
Consequently the walk's worst-case constant under Proposition~\ref{prop:walkbound} is
exactly $2$ in the limit, like the cover-directed scan's (Proposition~\ref{prop:chainfam}).
\end{proposition}

\begin{proof}
\emph{Optimum and face.} The $m$ diagonal constraints pair up the $2m$ coordinates, so
every feasible $y$ has $\sum y = \sum_i (y_{x_i} + y_{v_i}) \ge \sum_i d_i$; the point
$(0 \mid d)$ is feasible ($(x_i, v_{i+1})$ reads $0 + d_{i+1} \ge d_{i+1}$ and
$(x_{i+1}, v_i)$ reads $0 + d_i \ge s_i$), so $\opt = \sum_i d_i$ and every optimal $y$
is tight on every diagonal constraint: $y = (\varepsilon \mid d - \varepsilon)$ with
$\varepsilon_i := y_{x_i}$. Under this substitution $(x_i, v_{i+1})$ reads
$\varepsilon_i \ge \varepsilon_{i+1}$, $(x_{i+1}, v_i)$ reads
$\varepsilon_i - \varepsilon_{i+1} \le d_i - s_i = i$, the zero-weight constraints follow
from $y \ge 0$, and $y \ge 0$ is the box $0 \le \varepsilon_i \le d_i$, so $F$ is as
stated.

\emph{Pre-repair distance.} A face point $(\varepsilon \mid d - \varepsilon)$ agrees with
$\hat y = 0$ at $x_i$ exactly when $\varepsilon_i = 0$ and at $v_i$ exactly when
$\varepsilon_i = d_i$, never both, so it disagrees on at least $m$ coordinates;
$\varepsilon = 0$ attains $m$. Hence $\|y^* - \hat y\|_0 = m$.

\emph{Violated graph.} At $p = 0$ every positive-weight constraint is violated, with
residual equal to its weight. The diagonal is a perfect matching of $G(p)$ and either
side of the bipartition is a vertex cover, so $\tau(G(p)) = m$ by K\H{o}nig.

\emph{The run.} Start the walk at $x_1$. Its residuals are $d_1$ at $(x_1, v_1)$ and
$d_2 < d_1$ at $(x_1, v_2)$, so $\delta_{x_1} = d_1$ and the walk moves to the unique
attaining neighbour $v_1$. Inductively, suppose the walk arrives at $v_i$ ($i < m$)
having just deleted $x_i$, with $x_1, v_1, \dots, x_i$ deleted. The edges of $v_i$ were
$(x_i, v_i)$ and, for $i \ge 2$, $(x_{i-1}, v_i)$ (both endpoints already deleted)
and $(x_{i+1}, v_i)$, which remains: so $\delta_{v_i} = s_i$, and the walk moves to
$x_{i+1}$, deleting $v_i$ with that edge. At $x_{i+1}$ the remaining edges are
$(x_{i+1}, v_{i+1})$, of residual $d_{i+1}$, and (for $i + 1 < m$)
$(x_{i+1}, v_{i+2})$, of residual $d_{i+2} < d_{i+1}$: so $\delta_{x_{i+1}} = d_{i+1}$
with unique attaining neighbour $v_{i+1}$. At $v_m$ both incident edges
$(x_m, v_m)$ and $(x_{m-1}, v_m)$ are already deleted, so the pass ends with $v_m$ never
raised, and no edge of $G(p)$ remains ($(x_i, v_i)$ and $(x_i, v_{i+1})$
were deleted at $x_i$, and $(x_{i+1}, v_i)$ at $v_i$), so the walk terminates. No step
met a tie, so this is a resolution of the walk's free choices, and its output is
\[
  z \;=\; (d_1, \dots, d_m \mid s_1, \dots, s_{m-1}, 0).
\]

\emph{Post-repair distance.} A face point $(\varepsilon \mid d - \varepsilon)$ agrees
with $z$ at $x_i$ exactly when $\varepsilon_i = d_i$, at $v_i$ with $i < m$ exactly when
$\varepsilon_i = d_i - s_i = i$, and at $v_m$ exactly when $\varepsilon_m = d_m$. We
claim at most one index $i$ carries any agreement. Suppose indices $i < j$ both do, at
values $t_i \in \{d_i, i\}$ and $t_j \in \{d_j, j\}$ (for $j = m$ only $t_j = d_m$
occurs). Since $\varepsilon$ is non-increasing with steps
$\varepsilon_k - \varepsilon_{k+1} \le k$,
\[
  t_i \ge t_j
  \qquad\text{and}\qquad
  t_i - t_j \;\le\; \sum_{k=i}^{j-1} k .
\]
If $t_i = i$: both candidates for $t_j$ exceed $i$, since $j > i$ and
$d_j \ge d_m = m + 1 > i$, contradicting $t_i \ge t_j$. If $t_i = d_i$ and
$t_j = d_j$: telescoping, $d_i - d_j = \sum_{k=i}^{j-1} (k + 1) > \sum_{k=i}^{j-1} k$,
contradicting the step bound. If $t_i = d_i$ and $t_j = j$ with $j < m$: then
$d_i - j - \sum_{k=i}^{j-1} k = d_j + (j - i) - j = d_j - i \ge (m + 1) - i > 0$,
the same contradiction. So all agreements are at a single index; an index $i < m$
contributes at most one agreeing coordinate ($d_i \ne i$, as $d_i \ge m + 1 > i$), and
$i = m$ at most two, namely $x_m$ and $v_m$ together at $\varepsilon_m = d_m$. Hence
every face point disagrees with $z$ on at least $2m - 2$ coordinates, and
$\varepsilon \equiv d_m$ (non-increasing with zero steps, inside every box since
$d_m \le d_i$, and agreeing nowhere else since $d_m = m + 1$ avoids
$\{i : i < m\} \cup \{d_i : i < m\}$) attains $2m - 2$.

Both distances are minima over the whole face, computed against its exact description,
so the ratio is $(2m - 2)/m$ under this resolution, and by
Proposition~\ref{prop:walkbound} no resolution exceeds
$\|y^* - \hat y\|_0 + \tau(G(p)) = 2m$ on this family.
\end{proof}


\section{Experiments in full}\label{app:experiments}

\citeauthor{chen2022}'s bound is stated in the post-repair $\ell_0$ error, so that is what we
measure. Their proof counts the tight-subgraph deficiency of the warm-start dual and derives the $\ell_0$
bound from it, so $\ell_0$ is the parameter of their statement and the deficiency is the finer
quantity behind it. The solver we instrument is a standard successive-shortest-path (Hungarian)
solver and is not the algorithm they analyse, so what follows tests whether the post-repair $\ell_0$
error tracks warm-started work in a solver of this family.

Instances are $n \times n$ assignment problems of three kinds: uniform integer weights, clustered
(a few high-value blocks), and Euclidean ($w = C - \mathrm{dist}$). An optimal dual is computed by
LP, and $10\%$ of its coordinates are corrupted up or down to form the prediction.
Table~\ref{tab:exp} reports the mean post-repair $\ell_0$ error over ten draws per cell. In every
table the spread is the population standard deviation over the cell's draws (\texttt{numpy}'s
default, $\mathrm{ddof}=0$). Table~\ref{tab:measure} of Section~\ref{sec:exp} collects the $n = 200$ rows of
Tables~\ref{tab:exp} and~\ref{tab:solver}, together with the cover-directed scan's cells, which are
reported in the text below.

\begin{table}[ht]
\centering
\caption{Post-repair $\ell_0$ error by repair, against the perturbed reference optimum (an upper
bound on the face distance; the same reference for every repair). ``pre'' is the prediction's own
$\ell_0$ error, and the figure in parentheses is the increase over it. The walk uses
lowest-index start and tie-breaks. Mean ${\scriptstyle\pm}$\,std over ten draws per cell.}
\label{tab:exp}
\footnotesize
\setlength{\tabcolsep}{3pt}
\begin{tabular}{lrrrrrr}
\toprule
instance & $n$ & pre & sorted & forward & random & walk\\
\midrule
uniform   & 100 & $19.9\std{0.3}$ & $21.7\std{1.6}$ \,(+9\%) & $29.4\std{2.9}$ \,(+48\%) & $28.9\std{2.5}$ \,(+45\%) & $\mathbf{17.9}\std{1.3}$ \,($-$10\%)\\
uniform   & 200 & $39.8\std{0.6}$ & $43.4\std{2.3}$ \,(+9\%) & $54.9\std{4.9}$ \,(+38\%) & $54.7\std{5.9}$ \,(+37\%) & $\mathbf{32.9}\std{3.1}$ \,($-$17\%)\\
clustered & 100 & $19.9\std{0.3}$ & $22.1\std{2.3}$ \,(+11\%) & $28.1\std{2.5}$ \,(+41\%) & $28.7\std{2.7}$ \,(+44\%) & $\mathbf{17.9}\std{1.2}$ \,($-$10\%)\\
clustered & 200 & $39.7\std{0.6}$ & $47.9\std{3.6}$ \,(+21\%) & $67.2\std{6.2}$ \,(+69\%) & $64.2\std{8.5}$ \,(+62\%) & $\mathbf{35.3}\std{2.8}$ \,($-$11\%)\\
euclidean & 100 & $20.0\std{0.0}$ & $23.2\std{1.8}$ \,(+16\%) & $29.4\std{4.9}$ \,(+47\%) & $27.9\std{2.3}$ \,(+40\%) & $\mathbf{17.9}\std{1.4}$ \,($-$11\%)\\
euclidean & 200 & $39.8\std{0.4}$ & $48.2\std{3.2}$ \,(+21\%) & $63.8\std{6.3}$ \,(+60\%) & $66.2\std{5.3}$ \,(+66\%) & $\mathbf{35.5}\std{1.9}$ \,($-$11\%)\\
\bottomrule
\end{tabular}
\end{table}

Sorting increases the error less than the forward and random scans at both sizes. Its increase is $9$--$16\%$ against
$40$--$48\%$ at $n = 100$, and $9$--$21\%$ against $37$--$69\%$ at $n = 200$: a factor of $2.5$ to
$5.3$ and of $2.9$ to $4.2$ respectively. In coordinates the gap widens with $n$: forward exceeds
sorted by $6.0$--$7.7$ coordinates at $n = 100$ and by $11.5$--$19.3$ at $n = 200$, in every
family (the direction is stable across seeds, the exact ranges are not). As a ratio of increases the gap does
not widen: it narrows for five of the six combinations of family and scan, and the direction of
that change depends on the seed.

The sorted scan's post-repair error tracks the prediction's own $\ell_0$
error rather than the instance size, as Proposition~\ref{prop:sortedwarm} proves for $|A| = 1$.
Its sort costs $O(m \log m)$, negligible beside $O(m\sqrt{n})$.

The design is paired: every
repair sees the same instances and predictions. Among the scan orders, sorted has the smallest
mean in all twelve computed cells, the untabulated $n \in \{50, 150\}$ included, and beats forward
on $117$ of $120$ paired draws (one tie).

The walk has the lowest post-repair $\ell_0$ error in
every cell. It is below the prediction's own error in all twelve cell means, though not in every
draw ($108$ of $120$ strictly below, $9$ ties, $3$ above), and sorted never beats it in a paired
draw of this experiment. Its worst-case $\ell_0$ factor is~$2$ (Proposition~\ref{prop:walkbound}), the same
as the cover-directed scan's, and unlike the two-sided scans it is not $\ell_1$-non-expansive
(Section~\ref{sec:exp}).

On the same draws the cover-directed scan of
Theorem~\ref{thm:coverscan} ends within one coordinate of the prediction's own error in every
cell mean, between the sorted scan and the walk. At $n = 200$ it gives $39.5 \pm 0.8$,
$39.1 \pm 0.9$ and $39.7 \pm 0.5$ against $39.8$, $39.7$ and $39.8$ on the three families.
The same holds under a reversed or a random processing order of the cover. It fires $16$--$18$
times there while its bound allows $\tau(G(p)) = 19$--$21$, so the bound is loose on these
instances.

\subsection*{End to end: \texorpdfstring{$\ell_0$}{l0} and the solver's dual adjustments}

We test end to end whether the post-repair $\ell_0$ error predicts the work a warm-started solver
does. We start the solver from each repaired dual and count \emph{dual adjustments}, the quantity Corollary~\ref{cor:dinitz} bounds by
the post-repair gap. Alongside the four repairs of
Table~\ref{tab:exp} we run two members of the $\ell_\infty$-optimal class. One is the uniform
shift of Observation~\ref{obs:shift} (Appendix~\ref{app:delicate}). The other is the endpoint
shift, which raises only the endpoints of violated constraints by $\varepsilon/2$ and rounds those
up. It has the same $\ell_\infty$ displacement as the uniform shift on all $120$ draws but touches
far fewer coordinates, so the pair isolates the effect of the number of coordinates
touched. As a baseline we take the row-reduction start
$u_a = \max_b w_{ab}$, $v = 0$. Table~\ref{tab:solver} reports the counts.

\begin{table}[ht]
\centering
\caption{Dual adjustments of the warm-started solver, by starting point (mean
${\scriptstyle\pm}$\,std over ten draws per cell; instance families and corruption model as in
Table~\ref{tab:exp}). ``row-red.'' is the row-reduction start, and the walk is as in
Table~\ref{tab:exp}. The uniform shift (Observation~\ref{obs:shift}) and the endpoint shift have
the same realised $\ell_\infty$ displacement. Bold marks the cheapest start in a cell when a
repair achieves it; in the two uniform rows where the row-reduction start is cheapest nothing is
bolded.}
\label{tab:solver}
\footnotesize
\setlength{\tabcolsep}{2.5pt}
\begin{tabular}{lrrrrrrrr}
\toprule
instance & $n$ & row-red. & sorted & forward & random & walk & \begin{tabular}[b]{@{}r@{}}uniform\\ shift\end{tabular} & \begin{tabular}[b]{@{}r@{}}endpoint\\ shift\end{tabular}\\
\midrule
uniform & 50 & $39.9\std{5.6}$ & $28.9\std{6.8}$ & $29.8\std{7.4}$ & $29.8\std{6.4}$ & $\mathbf{17.9}\std{8.1}$ & $63.1\std{7.9}$ & $60.7\std{13.5}$ \\
uniform & 100 & $48.6\std{4.9}$ & $52.7\std{6.4}$ & $59.3\std{8.5}$ & $60.8\std{6.5}$ & $\mathbf{42.9}\std{5.6}$ & $126.9\std{5.8}$ & $141.8\std{25.8}$ \\
uniform & 150 & $45.8\std{3.3}$ & $62.3\std{9.1}$ & $76.1\std{11.5}$ & $78.2\std{13.6}$ & $53.2\std{7.8}$ & $179.9\std{7.4}$ & $218.3\std{30.2}$ \\
uniform & 200 & $42.7\std{3.4}$ & $72.6\std{8.6}$ & $84.7\std{10.2}$ & $83.1\std{11.7}$ & $56.3\std{10.4}$ & $253.1\std{10.1}$ & $252.7\std{27.3}$ \\
clustered & 50 & $55.8\std{13.4}$ & $24.8\std{6.0}$ & $25.3\std{5.4}$ & $25.7\std{4.4}$ & $\mathbf{17.0}\std{4.8}$ & $63.5\std{5.5}$ & $53.1\std{18.5}$ \\
clustered & 100 & $131.1\std{29.4}$ & $56.9\std{10.6}$ & $61.9\std{11.0}$ & $62.2\std{11.2}$ & $\mathbf{41.9}\std{8.7}$ & $130.8\std{10.4}$ & $156.5\std{23.0}$ \\
clustered & 150 & $247.8\std{42.7}$ & $84.0\std{14.6}$ & $96.9\std{12.7}$ & $92.3\std{11.8}$ & $\mathbf{63.2}\std{10.4}$ & $202.7\std{9.6}$ & $243.0\std{20.3}$ \\
clustered & 200 & $367.2\std{83.1}$ & $120.3\std{13.5}$ & $134.9\std{7.2}$ & $132.6\std{13.1}$ & $\mathbf{95.6}\std{10.7}$ & $281.4\std{11.9}$ & $309.7\std{24.1}$ \\
euclidean & 50 & $86.4\std{13.1}$ & $28.5\std{11.0}$ & $32.9\std{13.0}$ & $31.3\std{11.7}$ & $\mathbf{24.9}\std{9.5}$ & $70.4\std{8.1}$ & $59.8\std{18.9}$ \\
euclidean & 100 & $143.1\std{14.2}$ & $65.7\std{10.0}$ & $73.2\std{11.5}$ & $72.9\std{10.3}$ & $\mathbf{50.1}\std{7.9}$ & $135.3\std{10.9}$ & $143.8\std{21.5}$ \\
euclidean & 150 & $188.5\std{27.6}$ & $88.9\std{8.4}$ & $99.7\std{7.4}$ & $98.2\std{6.5}$ & $\mathbf{70.7}\std{11.3}$ & $211.7\std{10.5}$ & $245.3\std{21.6}$ \\
euclidean & 200 & $223.8\std{26.2}$ & $115.5\std{15.3}$ & $135.6\std{17.2}$ & $135.5\std{20.2}$ & $\mathbf{95.0}\std{10.5}$ & $265.6\std{10.6}$ & $295.1\std{27.1}$ \\
\bottomrule
\end{tabular}
\end{table}

\emph{First, across a set of five repairs, $\ell_0$ ranks them better than the alternatives, but
not inside one norm class.} The comparison is within an instance: we hold the instance fixed and
rank the repairs on it. The set consists of the sorted, forward and random scans, the walk and the uniform shift
(the row-reduction start, the cover-directed scan and the endpoint shift are not in it). Over this
set, the mean Spearman correlation with the adjustment count is $+0.95$ for the post-repair
$\ell_0$ error, against $+0.82$ for the $\ell_1$ error, $+0.80$ for the realised $\ell_1$
displacement and $+0.80$ for the dual gap. The last two coincide: every repair in the set only
raises coordinates, so the realised $\ell_1$ displacement is the gap plus a per-draw constant. The
margin is unlikely to be noise: the $\ell_0$ error wins $11$ of the $12$ cells, a one-sided sign-flip
permutation test over the per-cell differences gives $p = 0.0005$, and a bootstrap over instances
within the twelve fixed cells puts the difference between the $\ell_0$ and gap correlations at $+0.150$ with $95\%$ CI
$[+0.124, +0.174]$. Pooled over all $600$ points, the three quantities tie ($+0.9005$, $+0.9001$,
$+0.8998$). Pooling mixes in the size effect, which moves every repair together, so only the
within-instance comparison tells them apart. The margin also depends on the set: adding the
endpoint shift cuts $\ell_0$'s lead over the gap from $+0.150$ to $+0.085$, a $43\%$ reduction,
without reversing it. We therefore claim only that, among the quantities a repair can
guarantee, $\ell_0$ best predicts the solver's dual adjustments across this set of five repairs.

\emph{Second, both members of the $\ell_\infty$-optimal class that we run are expensive, and
$\ell_0$ does not rank inside it.} On uniform weights the uniform shift costs more than the
row-reduction start at every size, $5.9\times$ more at $n = 200$ ($253.1$ adjustments against
$42.7$). Across the twelve cells it exceeds the row-reduction start in seven. The endpoint shift also loses to the row-reduction start at every uniform size. But
the pair does not follow the first finding's ranking. The endpoint shift has strictly smaller post-repair $\ell_0$ error (between $0.34$ and
$0.76$ of the uniform shift's) and strictly smaller dual gap ($0.30$ to $0.72$) in all
twelve cells, yet it costs the solver more adjustments in eight of them ($141.8$ against $126.9$ at
uniform $n = 100$; $218.3$ against $179.9$ at $n = 150$). Neither $\ell_0$ nor the gap therefore
ranks these two correctly, and $\ell_\infty$-optimality is not the property that makes the
uniform shift expensive. We cannot explain this reversal, and it limits the first finding. We can
prove why both shifts are expensive: the cause is which edges are tight, and the size of
$\varepsilon$ does not enter.

\begin{proposition}[Each row without a tight edge costs an adjustment]\label{prop:tightrows}
Started from any feasible $(u \mid v)$, fractional or not, the solver of Table~\ref{tab:solver}
performs at least $U(u,v) = |\{a \in A : \min_b\, (u_a + v_b - w_{ab}) > 0\}|$ dual adjustments, one for every row that has no tight edge at the start. (In exact arithmetic; the implementation counts a step, and tests tightness, against a tolerance of $10^{-9}$, which replaces $0$ throughout.)
\end{proposition}
\begin{proof}
The solver is the successive-shortest-path method: phase $a$ grows an alternating tree of tight
edges from the unmatched row $a$, and when the tree is stuck it lowers $u$ on the tree's rows and
raises $v$ on its columns by the smallest positive slack leaving the tree, counting one adjustment.
So $u$ only decreases and $v$ only increases, and the tree of phase $a$ contains, besides $a$,
only rows matched in earlier phases; hence $u_a$ is untouched before phase $a$, and if $a$ had no
tight edge at the start then at the start of its phase every slack $u_a + v_b - w_{ab}$ is at
least its initial positive value, so the phase's first step is an adjustment. Distinct rows have
distinct phases.
\end{proof}

\begin{corollary}\label{cor:shiftrows}
For integral duals, the uniform shift with an odd largest violation $\varepsilon$ raises every
slack by $\varepsilon + 1 > \varepsilon$, so it leaves no tight edge and forces at least $|A|$
adjustments, however small $\varepsilon$ is; with $\varepsilon$ even every slack rises by exactly
$\varepsilon$, and the only tight edges left are the maximally violated constraints. With $\varepsilon$ odd the endpoint shift leaves every touched row without a tight edge and forces at least their number; with $\varepsilon$ even, at the touched rows it too leaves tight exactly the maximally violated constraints, so it forces at least the number of touched rows not incident to one. The row-reduction start has $U = 0$.
\end{corollary}
\begin{proof}
Let $(u \mid v)$ be the clamped integral prediction. Since $w$ is integral, every slack
$u_a + v_b - w_{ab}$ is an integer of at least $-\varepsilon$, with equality exactly on the
maximally violated constraints. The uniform shift adds $\lceil\varepsilon/2\rceil$ to every
coordinate, so it adds $2\lceil\varepsilon/2\rceil$ to every slack: $\varepsilon + 1$ for odd
$\varepsilon$ and $\varepsilon$ for even $\varepsilon$. The endpoint shift adds the same
$2\lceil\varepsilon/2\rceil$ to an edge with both endpoints touched. An edge from a touched row to
an untouched column was not violated, since both endpoints of a violated constraint are touched, so
its slack was non-negative and becomes at least $\lceil\varepsilon/2\rceil \ge 1$. So for odd
$\varepsilon$ no touched row keeps a tight edge, and for even $\varepsilon$ the tight edges at a
touched row are exactly the maximally violated constraints incident to it. The row-reduction start
makes a heaviest edge of every row tight. Proposition~\ref{prop:tightrows} turns each row without a
tight edge into one adjustment.
\end{proof}

On the $120$ draws of Table~\ref{tab:solver}, $\varepsilon$ is odd in $62$, and there the
uniform shift has $U = |A|$ exactly. $U$ accounts for $70$--$83\%$ of its adjustments in every
cell, against $22$--$45\%$ for the scans and the walk. The endpoint shift's $U$ runs from
$0.41|A|$ to $0.93|A|$ in cell means ($0.12|A|$ to $1.00|A|$ across individual draws). The
count $U$ explains the first half of the second finding: with odd $\varepsilon$ the uniform shift
needs one adjustment per row before it does anything else, while the row-reduction start, which
leaves a tight edge in every row, needs none. It does not explain the second half: $U$ is smaller for
the endpoint shift than for the uniform shift in $119$ of the $120$ draws, so the endpoint shift's
extra adjustments come from the later steps of the phases. As a ranking, $U$ is nearly as good as
$\ell_0$: within an instance, over the same five repairs, it correlates with the adjustment count
at $+0.93$, against $\ell_0$'s $+0.95$ and the gap's $+0.80$.\footnote{One consequence of Proposition~\ref{prop:tightrows} separates the gap from the count. From any optimal dual the solver adjusts nothing, since its tight graph contains a perfect matching. From $(u^* + t \mid v^* + t)$ it adjusts exactly $|A|$ times, once per row, for every $t > 0$ above the implementation's tolerance, while the gap is $2|A|t$. For integral duals the count is bounded by the gap, by the argument of Corollary~\ref{cor:dinitz}: each adjustment lowers $\sum_a u_a + \sum_b v_b$ by an integer $\delta \ge 1$ while the dual value stays at least $\opt$. For fractional starts, in exact arithmetic, no function of the gap or of any $\ell_p$ displacement bounds the count. Only the dimension does, trivially: at most one adjustment per step of a phase, so at most $|A||B|$. (The implemented counter, with its tolerance $\theta = 10^{-9}$, satisfies $\mathrm{adj} \le g/\theta$, so its discontinuity is at the scale of $\theta$.)}

\emph{Third, a $10\%$-corrupted optimum is not always a better start than the row-reduction
start.} On uniform weights the row-reduction start beats every scan-order repair from
$n = 100$ on, and every repair, the walk included, at $n \in \{150, 200\}$. Warm-starting helps
where the row-reduction start is weak (clustered: $367.2$ adjustments from it against $95.6$ for
the best repair at $n = 200$, a difference far outside the $\pm 83.1$ of that cell, the largest spread in
the table). The cold start of
Proposition~\ref{prop:sortedcold} (repair the zero prediction) is a structural result and we do
not recommend it: as a solver start it loses to the row-reduction start by $1.4$--$7.3\times$ across
these families. The walk is again the cheapest repair start in every cell, now on solver work. The
cover-directed scan is the repair with both a guarantee
and a measured $\ell_0$ increase near zero. As a start on the same draws it costs
$66.6 \pm 9.7$, $110.4 \pm 11.0$ and $110.2 \pm 15.0$ adjustments at $n = 200$ on the three
families, against the sorted scan's $72.6$, $120.3$ and $115.5$ and the walk's $56.3$, $95.6$ and
$95.0$. It lies between the two in all twelve cells, beating the sorted scan on $95$
of $120$ draws and the walk on $4$.

Among the scan orders, the $\ell_0$ differences of Table~\ref{tab:exp} do carry through to solver
work. Sorted has the smallest mean in all twelve cells of Table~\ref{tab:solver} as well, and beats
forward on $104$ of $120$ paired draws ($5$ ties), with per-cell mean margins from $0.5$ to $20.1$
adjustments. The exact one-sided sign test on the $115$ nonzero differences gives
$p = 1.9\times10^{-20}$ (the signed-rank test agrees, but its normal approximation is unreliable
this far in the tail). This is a comparison among scan orders only, and the walk is not one of them.

\paragraph{The algorithm the bound is about.} \citeauthor{chen2022}'s solver (their Algorithm~1)
is itself of the successive-shortest-path family: a maximum-cardinality matching on the tight subgraph of $\hat y'$, then repeated
phases of Dijkstra and maximum flow. We implemented it line by line and ran it on the same $120$ draws and
the same repaired duals.

One quantity we measure is the deficiency $n - |M|$ of the tight subgraph after its line~4, the quantity their
proof counts. It is the flow the algorithm must still push, and the maximality of $M$ determines it
from $\hat y'$ alone. Over the same five repairs, post-repair $\ell_0$ ranks it at $+0.98$, against
$+0.82$ for the $\ell_1$ error and $+0.81$ for the gap, with $\ell_0$ strictly best in all twelve
cells. Of the repair pairs that $\ell_0$ orders strictly, only $0.9\%$ come out reversed. As with
the adjustments, the result depends on the set of repairs: adding the endpoint shift raises all three ($+0.99$, $+0.90$,
$+0.89$) and cuts $\ell_0$'s lead over the gap from $+0.177$ to $+0.100$, without reversing it.

The other is the number of while-loop iterations, which their $(m + n\log n)$ factor multiplies.
Their Lemma~B.9 bounds it by $\ell_0$, but it is a phase count, with a median well below $\ell_0$
here. $\ell_0$ does not rank this count: it scores $+0.58$, against $+0.61$ for the $\ell_1$ error
and $+0.60$ for the gap, with $17.6\%$ of the strictly $\ell_0$-ordered pairs reversed and a
further $18.6\%$ tied. The pseudocode also leaves this count open: line~12 asks for a ``shortest
path distance in $G_f$'' without saying whether the residual arcs out of $t$ enter the Dijkstra,
and line~16 does not say which maximum flow it returns. The count therefore depends on the
implementation, unlike the deficiency. Under the other reading $\ell_0$ does not lead either; we
quote no single figure for it. On total counted work, which on these instances is dominated by the
Dijkstra term rather than the Ford--Fulkerson one, $\ell_0$ again does not lead ($+0.47$ against
$+0.54$).

\subsection*{The regime where Theorem~\ref{thm:sharpwarm} has content}

Theorem~\ref{thm:sharpwarm} has content only when
$\|y^* - \hat y\|_0 \le \bigl||B| - |A|\bigr|$, so Table~\ref{tab:rect} measures rectangular
instances with $\min(|A|,|B|) \ll \max(|A|,|B|)$. All $80$ draws are well inside that regime
($\|y^* - \hat y\|_0$ runs from $13$ to $37$, against $\bigl||B|-|A|\bigr|$ between $280$ and
$597$), so the theorem has content on the whole table. It is also loose there: it allows an
increase of $2\min(|A|,|B|)-1$, between $5$ and $39$, where the sorted scan has none.

In cell means every scan lies within one standard deviation of the
prediction's own error and of the others. Cell means are the wrong statistic here, because the
design is paired. All four
repairs see the same instance and the same corrupted prediction, and
the spread within a column comes from the draws: the sorted column's spread equals
that of pre in all eight cells (the table does not print it), because the two agree draw for draw.
Of the $80$ draws, $49$ are no-ops: on these very unbalanced
shapes a $10\%$ corruption often leaves the prediction already feasible, and then every repair is
the identity. They range from $2$ of $10$ at $20\times300$ to $10$ of $10$ on clustered
$3\times600$, whose row is therefore degenerate, and they add a tie to every comparison
and nothing else. On the remaining $31$ draws the orders separate with no exception in
either direction. The sorted scan ends at exactly $\|y^*-\hat y\|_0$ on all $31$: an increase
of $0$, against the $2\min(|A|,|B|)-1$ the theorem allows. The forward scan is strictly
worse on $22$ of the $31$ and never better, the random scan on $25$ and never better. The walk is
strictly below the prediction's own error on $30$ and never above it, so it beats the sorted scan
on $30$ draws and loses on none, the same ordering as in the solver table. Since every nonzero
difference has the same sign, the exact one-sided signed-rank test coincides with the sign test and
returns $2^{-22}$, $2^{-25}$ and $2^{-30}$ respectively.

The sorted scan's equality does not come from inaction: on all $31$ draws it moves between $2$ and
$6$ coordinates, and its wrong set differs from the prediction's on every one, making $42$ wrong
coordinates correct and $42$ correct coordinates wrong across the $31$. But all these $\ell_0$
figures are distances to the fixed perturbed reference optimum, whereas
Theorem~\ref{thm:sharpwarm} measures distance to the optimal face, and on $21$ of the $31$ draws that face is not a single
point. The equality is therefore evidence for additivity with constant $0$ in this regime; it does not
measure the face distances that the theorem bounds.
The absence of exceptions is a property of this pinned stream:
re-running the same design at other seeds gives occasional exceptions in every direction (at seed
$101$ the sorted scan exceeds $\|y^*-\hat y\|_0$ on $2$ of $32$ non-no-op draws and the walk loses
one; at seed $303$ the walk loses two; at seed $404$ the sorted scan loses one draw to the forward
scan). Across seeds the direction and the margins are stable, but the absence of exceptions is not.
Order-induced increases in $\ell_0$ error happen on balanced instances, where Theorem~\ref{thm:balancedlb}
shows that the repair itself can cause them.

The uniform shift, left out of Table~\ref{tab:rect}, has mean
post-repair $\ell_0$ error between $89$ and $260$ on the seven rows where any draw was
infeasible, against at most $33$ for every scan, with standard deviations from $\pm121$ to $\pm221$.
This reflects the dependence of its single global $\varepsilon$ on the worst violated constraint.
At this scale Lemma~\ref{lem:budget} is vacuous, so we claim no face distance for the uniform shift.

\begin{table}[ht]
\centering
\caption{Post-repair $\ell_0$ error on unbalanced instances, where Theorem~\ref{thm:sharpwarm} has
content (mean ${\scriptstyle\pm}$\,std over ten draws per cell; ``walk'' as in
Table~\ref{tab:exp}; the uniform shift is discussed in the text). ``no-ops'' counts the draws of
each row whose corrupted prediction was already feasible, so that every repair is the identity
($49$ of the $80$ draws). The cell means therefore look flat; read the table per draw (see
text).}
\label{tab:rect}
\footnotesize
\setlength{\tabcolsep}{2.5pt}
\begin{tabular}{lrrrrrrrrr}
\toprule
instance & shape & $|A|{+}|B|$ & $2\min{-}1$ & no-ops & pre & sorted & forward & random & walk\\
\midrule
uniform & $5\times500$ & $505$ & $9$ & $8/10$ & $24.6$ & $24.6\std{3.6}$ & $24.9\std{3.8}$ & $25.2\std{4.0}$ & $24.3\std{3.6}$ \\
uniform & $10\times400$ & $410$ & $19$ & $6/10$ & $22.0$ & $22.0\std{2.9}$ & $23.2\std{3.4}$ & $23.5\std{3.4}$ & $21.3\std{3.2}$ \\
uniform & $20\times300$ & $320$ & $39$ & $2/10$ & $17.2$ & $17.2\std{2.3}$ & $19.0\std{3.1}$ & $19.0\std{3.3}$ & $16.3\std{2.6}$ \\
uniform & $3\times600$ & $603$ & $5$ & $9/10$ & $32.2$ & $32.2\std{3.0}$ & $32.6\std{3.4}$ & $32.2\std{3.0}$ & $32.1\std{2.9}$ \\
clustered & $5\times500$ & $505$ & $9$ & $7/10$ & $22.7$ & $22.7\std{3.0}$ & $23.2\std{3.2}$ & $23.1\std{3.4}$ & $22.3\std{3.0}$ \\
clustered & $10\times400$ & $410$ & $19$ & $5/10$ & $21.6$ & $21.6\std{2.4}$ & $22.4\std{2.3}$ & $22.5\std{2.5}$ & $21.1\std{2.5}$ \\
clustered & $20\times300$ & $320$ & $39$ & $2/10$ & $16.9$ & $16.9\std{2.9}$ & $18.2\std{3.5}$ & $18.8\std{3.8}$ & $16.0\std{2.8}$ \\
clustered & $3\times600$ & $603$ & $5$ & $10/10$ & $32.3$ & $32.3\std{2.5}$ & $32.3\std{2.5}$ & $32.3\std{2.5}$ & $32.3\std{2.5}$ \\
\bottomrule
\end{tabular}
\end{table}

\subsection*{Drifted predictions}

Corrupting a fraction of the true optimum is the standard experiment in this line, but a
prediction in practice is more likely to be the dual of an earlier instance. Table~\ref{tab:drift} re-runs the solver
experiment with predictions obtained by solving one instance, moving every weight by a random
integer between $-10$ and $10$ and clipping the result at $0$ (the bound $10$ is about a tenth of the
range of the uniform and Euclidean weights), and predicting the new dual by the old one. Under drift the $\ell_0$ error saturates
($114.7$ of $120$ coordinates at $n = 60$): drift perturbs every dual coordinate a little, which is
the error structure the $\ell_0$ parameterisation has nothing to say about. This limits every
$\ell_0$-parameterised bound in this line, ours included. The ranking of
the repairs also compresses (the walk is still first among them in every cell), while the
row-reduction start wins four of the six cells, three of them by more than one standard deviation.
The fourth, euclidean $n = 60$, is inside the noise ($79.2 \pm 13.6$ against the walk's
$82.4 \pm 7.9$), and the row-reduction column remains the most variable. The scan-order ranking of
Table~\ref{tab:solver} does not hold under drift either: the sorted scan has the largest mean of the
three scans in five of the six cells and the smallest in none.

\makeatletter\setlength{\@fptop}{0pt}\makeatother
\begin{table}[ht]
\centering
\caption{Dual adjustments under drifted predictions: the previous instance's optimal dual
predicts the drifted instance's (mean ${\scriptstyle\pm}$\,std over ten draws per cell).
Columns are as in Table~\ref{tab:solver} without the endpoint shift, plus ``pre'', the prediction's
own $\ell_0$ error, near-saturated at every size.}
\label{tab:drift}
\footnotesize
\setlength{\tabcolsep}{3pt}
\begin{tabular}{lrrrrrrrr}
\toprule
instance & $n$ & pre & row-red. & sorted & forward & random & walk & \begin{tabular}[b]{@{}r@{}}uniform\\ shift\end{tabular}\\
\midrule
uniform & 60 & $114.7$ & $65.3\std{6.8}$ & $95.1\std{4.4}$ & $95.5\std{4.5}$ & $94.8\std{4.6}$ & $77.1\std{7.4}$ & $103.3\std{6.6}$ \\
uniform & 120 & $226.0$ & $101.4\std{8.1}$ & $183.8\std{7.4}$ & $182.6\std{5.2}$ & $177.7\std{5.5}$ & $164.1\std{5.8}$ & $185.5\std{6.4}$ \\
clustered & 60 & $116.1$ & $104.3\std{22.8}$ & $99.1\std{6.0}$ & $96.3\std{6.4}$ & $94.8\std{6.2}$ & $83.5\std{7.4}$ & $101.5\std{6.3}$ \\
clustered & 120 & $230.6$ & $277.4\std{56.8}$ & $189.8\std{6.3}$ & $185.1\std{6.2}$ & $180.3\std{7.0}$ & $164.8\std{10.0}$ & $191.5\std{8.4}$ \\
euclidean & 60 & $114.5$ & $79.2\std{13.6}$ & $98.8\std{5.8}$ & $95.7\std{5.9}$ & $91.6\std{5.9}$ & $82.4\std{7.9}$ & $102.9\std{4.9}$ \\
euclidean & 120 & $226.4$ & $146.0\std{18.6}$ & $195.4\std{8.0}$ & $192.6\std{6.7}$ & $184.0\std{6.8}$ & $170.2\std{8.6}$ & $190.9\std{3.4}$ \\
\bottomrule
\end{tabular}
\end{table}


\clearpage
\section{Reproducibility}\label{app:repro}

\paragraph{Code, licence and environment.} Every theorem check, every table and every displayed
constant in this paper is regenerated by the code at \coderepo, released under the \codelicence{} licence. It
does this with seven Python programs. Each program runs as
\texttt{python3 <file>} with no arguments, and every configuration and seed is pinned in the file.
\texttt{verify\_all.py} needs only the Python standard library; every other program needs NumPy
\citep{harris2020} and SciPy \citep{virtanen2020}. The code was tested with Python~3.10,
NumPy~1.26 and SciPy~1.15, on an Apple M3 Pro laptop (12 cores, 36\,GB RAM, macOS), CPU only, and
the times below refer to that machine.

\paragraph{The programs.}
\begin{itemize}
\item \texttt{verify\_all.py} (standard library only; about five minutes): every exhaustive
check, namely non-expansion in each norm (Theorems~\ref{thm:nonexp} and~\ref{thm:pack},
Corollary~\ref{cor:lip}), the two failing repairs of Proposition~\ref{prop:fail}, the
order-dependence table (Table~\ref{tab:order}) including its $\ell_\infty$ and $\ell_\infty^\pm$
columns, the cold-start sweep, the sorted-scan warm-start sweep, the sharpness family of
Proposition~\ref{prop:sharpattained}, and the witness and sweep of
Section~\ref{sec:exp} that compare the walk with the sorted scan.
\item \texttt{experiments.py} (about $20$ seconds): Table~\ref{tab:exp}, from a pinned
configuration.
\item \texttt{solver.py} (about $3$ seconds): the warm-startable successive-shortest-path solver
of Appendix~\ref{app:experiments} and its correctness suite.
\item \texttt{experiments\_itcs.py} (about two minutes): Tables~\ref{tab:solver},
\ref{tab:rect} and~\ref{tab:drift}, the correlation figures and their significance tests, the
paired comparison of the sorted and forward scans with its exact sign test, the per-draw
analyses of Table~\ref{tab:rect} with the alternate-seed runs, the cold-start comparison,
and the tight-row count $U$ of Proposition~\ref{prop:tightrows}.
\item \texttt{verify\_balanced.py} (about one minute): Theorem~\ref{thm:balancedlb} on
$\mathcal{B}_n$, the $\ell_\infty$ and $\ell_\infty^\pm$ reversal on the same family, and the
$\ell_1$ distance to the optimal face, equal for both orders.
\item \texttt{chen\_alg1.py} with \texttt{measure\_chen.py} (about one minute): a line-by-line
implementation of \citeauthor{chen2022}'s solver (their Algorithm~1 in \texttt{arXiv:2204.12055v1}) and the
deficiency and iteration-count measurements on the draws of Table~\ref{tab:solver}.
\item \texttt{verify\_dichotomy.py} (about twelve minutes; imports \texttt{solver.py},
\texttt{experiments.py} and \texttt{experiments\_itcs.py}): the results of
Section~\ref{sec:composition} and Appendix~\ref{app:dichotomy}, Lemma~\ref{lem:balancedface}, the
cold-start results of Appendix~\ref{app:sorted}, Proposition~\ref{prop:tightrows}, and the
cover-directed cells of Table~\ref{tab:measure}; details below.
\end{itemize}
The seven programs together take about twenty minutes. The repository also holds exploratory
scripts and checks that this paper does not use.

\paragraph{Table labels in the programs' output.} The repository's README maps the labels in the
programs' output to this paper's numbering. Table~\ref{tab:measure} of Section~\ref{sec:exp}
collects the $n = 200$ rows of Tables~\ref{tab:exp} and~\ref{tab:solver}; its cover-directed cells
are printed by \texttt{verify\_dichotomy.py}, which imports both pinned streams.

\paragraph{Validation.} \texttt{solver.py} is verified against an exact reference
(SciPy's \texttt{linear\_sum\_\allowbreak{}assignment}) on $1{,}954$ instances, all of which match
exactly. They come from $333$ weight matrices spanning dense, heavily degenerate (all-equal,
$\{0,1\}$, small-alphabet, constant-row, rank-one and large tie-block) and sparse, near-disconnected
families, each solved from several feasible starting duals, $1{,}293$ of them warm. The suite also
checks invariants of the adjustment counter as well as its output. On all $1{,}954$ runs the solver
performs exactly $n$ augmenting phases, and, since the data are integral, the reported number of
dual adjustments never exceeds the initial dual gap $\sum_i u^0_i + \sum_j v^0_j - \opt$, as
Corollary~\ref{cor:dinitz} requires. (The solver is a square-matrix routine, so the rectangular
experiment of Table~\ref{tab:rect} reports no solver work.)
The implementation of \citeauthor{chen2022}'s solver performs $1{,}560$ solves on the stream
of Table~\ref{tab:solver}. It returns an optimal matching on each of the $1{,}440$ whose
weight is checked against \texttt{scipy.optimize.linear\_sum\_assignment}; the $120$ solves from
the row-reduction start are run but not weight-checked. At runtime it is checked against that
paper's Lemma~B.5 at the base case and at every matched edge, and against Lemmas~B.6, B.8 and B.9
throughout. The exhaustive per-iteration B.5 sweep is available under \texttt{check=True} and is
off in the pinned run.

The non-expansion checks all return zero violations with maximum ratio exactly 1.000000.
They cover $7.7\times10^5$ comparisons per norm for $p \in \{1, 1.5, 2, 3, \infty\}$
(Theorem~\ref{thm:nonexp}), $1.0\times10^6$ bipartite and $9.8\times10^5$ non-bipartite comparisons
for the covering case, and $8.2\times10^5$ for the packing mirror (Theorem~\ref{thm:pack}).
Corollary~\ref{cor:lip}'s two-input Lipschitz property is checked separately, on $36{,}000$
comparisons ($6{,}000$ per norm for $p \in \{1,1.5,2,3,7,\infty\}$ over $2{,}974$ bipartite and
$3{,}026$ non-bipartite systems, predictions allowed negative entries): again zero violations,
maximum ratio exactly $1.000000$. For Proposition~\ref{prop:fail} the same search reports
$155{,}270$ violations for the alternating projection and $37{,}674$ for the integral split
variant, with ratios reaching $3$ on those small instances. The two explicit families of its proof
are checked through $k = 12$ and drive the ratio to $|V| - 1$ and $|V|$ respectively, as the proof
predicts.

\texttt{verify\_balanced.py} checks Theorem~\ref{thm:balancedlb} on $\mathcal{B}_n$ for
$n = 2, \dots, 8$: $\opt = n^2$, pre-repair face distance $1$, and post-repair face distance $n$
with $n$ firings under every weakly-descending order. The orders are enumerated exhaustively for
$n \le 4$ ($20{,}736$ orders at $n = 4$) and sampled otherwise; the order-independence itself is
proved in Step~3 of the theorem. It also checks that the increasing scan ends at distance $1$,
while the dual gap, the realised $\ell_1$ displacement and the $\ell_1$ distance to the optimal
face (equal to $n$) are identical for the two orders. Its face distances are also exact minima over
the fractional optimal face, computed as a mixed-integer program.

\texttt{verify\_dichotomy.py} runs in numbered sections, each of which ends by reporting that all
its assertions hold. It checks the following.
\begin{itemize}
\item Theorem~\ref{thm:noadditiveclass}. It certifies the gadget: singleton face by the LP extremes
of every coordinate; the four agreement linear programs infeasible, so the conclusion holds over
real outputs; for two copies, exactly one admissible integral point in the box, at distance $8$
from $6$, reached by every scan order and by all $34{,}560$ weakly-descending tie-breaks. It runs
the exhaustive search of Remark~\ref{rem:noadditivescope}: $201$ witnesses on $1 \times 3$
systems, none on $1 \times 2$, every one saturating with ratio $4/3$ and confirmed over real
outputs by linear programming, the gadget among the three minimal ones.
\item Theorem~\ref{thm:coverscan}. The cover-directed scan runs on $3{,}000$ random bipartite
instances over every minimum cover, every processing order and every tie among largest residuals
($6{,}082$ outputs, all checked against the sharp bound $\|y^* - \hat y\|_0 + |S|$; the maximum
ratio on this sample is $4/3$). The maximal-matching variant runs on the same instances and on
$500$ non-bipartite systems (face distances there by half-integral enumeration, cross-checked
against linear programs over the fractional face; maximum ratio $2$ on that sample). A hill-climb
witness attains the additive term at ratio $7/4$ (section [2b]).
\item Proposition~\ref{prop:walkbound}. The walk runs over all $10{,}270$ tie-break resolutions of
$3{,}000$ instances (maximum ratio $5/3$ on that sweep), on its family and two copies of it under
every resolution, and on the $3 \times 4$ witness of Section~\ref{sec:exp}, with its reachable set
$\{2,3,4,6\}$.
\item Proposition~\ref{prop:ladderfam}. The ladder family is verified through $m = 7$ ($\opt$, the
pre-repair distance $m$, $\tau = m$, and the designed run's distance $2m-2$, all as exact minima
over the fractional face), with every walk resolution enumerated for $m \le 6$ and the worst always
exactly $2m-2$ (section [14]).
\item Lemma~\ref{lem:balancedface}. The face of $\mathcal{B}_n$ is verified for $n \le 12$ by
linear programming, with the $\ell_\infty$ and $\ell_\infty^\pm$ values of
Remark~\ref{rem:balancedorder} and all weakly-descending and weakly-increasing tie-breaks for
$n \le 6$.
\item Proposition~\ref{prop:tightrows} and Corollary~\ref{cor:shiftrows}: the tight-row bound on
$360$ starts and the parity of $\varepsilon$ on $120$ instances.
\item Table~\ref{tab:exp}. Importing \texttt{experiments.py}'s pinned stream, it reproduces the six
tabulated rows to the printed digit, alongside the cover-directed scan's cells.
\item Theorem~\ref{thm:composed}. The packing mirror of the cover-directed scan is checked over
every cover, order and tie on $200$ random square instances, against the unbounded fractional face.
\item Table~\ref{tab:solver}. From \texttt{experiments\_itcs.py}'s pinned stream, it reproduces
the row-reduction, sorted and walk cells at $n \in \{100, 200\}$ to the printed digit, alongside the
cover-directed start's adjustments.
\item Theorem~\ref{thm:threehalves}. The family of (b) is recomputed from the $2^k$ points $y_x$ of its proof
alone for $k \le 4$. The bound $r + \lfloor r/2 \rfloor$ of (a) is checked against the exact
minimum over admissible outputs on every $1 \times 2$ and $1 \times 3$ instance with weights and
predictions in $\{0, \dots, 4\}$ and on random near-face instances, together with the two-round
construction itself. The exact minimum is computed by linear programming over the polytope
$\{z \in \dfeas : e(z) \le 0\}$, with $e$ expanded over all sign vectors.
\item Proposition~\ref{prop:faceadditive}: its factor~$2$ is exercised numerically (section [7]).
\item Proposition~\ref{prop:chainfam}: the chain family for $k \le 8$.
\item Theorem~\ref{thm:coldfivehalves}. It checks $\lfloor 5\min(|A|,|B|)/2 \rfloor$ against every
tie-break's cold output on $300$ random instances and the $4 \times 5$ witness, with exact face distances
(section [17]). For the further parts it checks the families of
Proposition~\ref{prop:coldfamilies}, with LP-exact face distances, unique optimal duals and
order-free cold outputs; the six cold outputs of the $6 \times 8$ witness at face distances $11$ to
$14$, by an exact mixed-integer program validated against integral enumeration; and the bound
$\lfloor 5\min(|A|,|B|)/2 \rfloor - 1$ on $355$ cold outputs over every tie-break (section [20]).
\item Proposition~\ref{prop:coldrefuted}: the witness and its direct-sum growth (section [12]).
\end{itemize}
Most of its running time goes to Theorems~\ref{thm:threehalves} and~\ref{thm:coldfivehalves} and the ladder family.

\paragraph{Exactness and order-robustness.} All $\ell_0$, $\ell_1$ and $\ell_\infty$ figures in
Table~\ref{tab:order}, and the cold-start sweep, are computed exactly as minima over the full
optimal face. The $\ell_0$ and $\ell_1$ figures come from enumerating the integral face, which the
total-unimodularity argument of Appendix~\ref{app:early} justifies for those two norms only. The
$\ell_\infty$ and $\ell_\infty^\pm$ figures come from exact rational minimisation over the full
continuous face (the text derives the same two values in closed form). The corresponding figures
for $\mathcal{B}_n$ are derived in closed form in Remark~\ref{rem:balancedorder} and recomputed by
\texttt{verify\_balanced.py} by linear programming over the same full face. The witness of Section~\ref{sec:exp}
that compares the walk with the sorted scan, and its $400$-instance sweep, are also recomputed in
\texttt{verify\_all.py}, with exact face distances. The witnesses of Appendix~\ref{app:sorted} are
verified exactly under every weakly-descending tie-break order (zero-weight pairs come last in
any such order and provably never fire, so enumerating the permutations of the positive weights is
exhaustive), as is the sharpness family of Definition~\ref{def:sharpfam} through
$\min(|A|,|B|) = 12$. Where we claim order-robustness the output is a single point, the sharpness
family and the $3\times4$ cold witness included. The one exception is the $2\times6$ warm witness,
which attains its increase of $3$ under $48$ of its $144$ weakly-descending orders and $1$ under the rest;
for it we claim only that the increase is reached. A randomised sweep
($500$ instances per envelope, shapes up to $4 \times 4$) reports maximum slack $-1$ and $-2$
against $2\min(|A|,|B|)$ and Theorem~\ref{thm:sharpwarm}; the first envelope is the one
Proposition~\ref{prop:coldrefuted} breaks at $4 \times 5$. \texttt{verify\_all.py}'s two
independent face routines (exact minimisation over the fractional face and brute-force enumeration
of the integral face) agree on $9{,}000$ (instance, point) pairs, and its matching routine agrees
with brute force on $3{,}000$ instances.

\paragraph{Main commands.}
\begin{verbatim}
python3 verify_all.py       # no violation, largest ratio exactly 1.000000
python3 verify_balanced.py  # distance 1 becomes n, every descending order
python3 verify_dichotomy.py # sections end 'all assertions in [k] hold'
python3 measure_chen.py     # Chen et al.'s Algorithm 1 on the solver draws
python3 solver.py           # all 1954 instances match the SciPy optimum
python3 experiments.py      # the post-repair l0 table
python3 experiments_itcs.py # solver, unbalanced, drift tables, statistics
\end{verbatim}
